% ============================================================================
% Fichier  : partie1/P01_C02_menaces_attaques.tex
% Titre    : Menaces, Vulnérabilités et Attaques
% Auteur   : MINKA MI NGUIDJOÏ Thierry Emmanuel
% Version  : 1.0 -- novembre 2026
% ============================================================================
% Base : brouillon « Livre I », section « Bases de la cybersécurité ».
% Sources : chaîne d'attaque (Hutchins, Cloppert, Amin, 2011) ; MITRE ATT&CK
% (Strom et al., 2018) ; CVE ; CVSS v3.1 (FIRST, 2019) ; loi n° 2010/012,
% art. 65 (2), 66, 69, 70, 86 (texte officiel, cités au Chapitre 11).
% ============================================================================

\chapter{Menaces, Vulnérabilités et Attaques}
\label{chap:menaces}
\chaptermark{Menaces et attaques}

\epigraph{%
    « L'attaquant n'a besoin de réussir qu'une fois~;\\
    le défenseur doit tenir à chaque fois, sauf s'il sait détecter. »%
}{Principe du cours}

\niveauChapitre{Fondamental}{Chapitre~\ref{chap:proprietes}
(propriétés, menace, vulnérabilité).}

\begin{boxobjectifs}
\begin{itemize}
    \item Classer les sources de menace selon leur motivation et leurs
          capacités.
    \item Décrire les principales familles de logiciels malveillants.
    \item Expliquer l'ingénierie sociale et l'hameçonnage.
    \item Décrire une attaque comme une chaîne, et la lire avec MITRE
          ATT\&CK.
    \item Lire une fiche CVE et un score CVSS.
\end{itemize}
\end{boxobjectifs}

\bigskip

% ============================================================================
\section*{Question directrice : qui attaque, comment, et pourquoi a-t-il réussi~?}
% ============================================================================

L'attaquant de Garoua n'a utilisé aucune technique nouvelle. Un courriel
trompeur, un mot de passe volé, un correctif manquant, des sauvegardes à
portée de main~: tout était connu, documenté, et évitable. Comprendre
comment les attaquants travaillent, c'est savoir où les attendre.

\medskip
\textbf{Plan du chapitre.}
\begin{enumerate}
    \item Les sources de menace (section~\ref{sec:c02-sources}).
    \item Les logiciels malveillants (section~\ref{sec:c02-malware}).
    \item L'ingénierie sociale (section~\ref{sec:c02-social}).
    \item Les attaques contre les réseaux et les services (section~\ref{sec:c02-reseau}).
    \item La chaîne d'attaque et MITRE ATT\&CK (section~\ref{sec:c02-chaine}).
    \item Vulnérabilités~: CVE et CVSS (section~\ref{sec:c02-cve}).
\end{enumerate}

\separateur

% ============================================================================
\section{Les Sources de Menace}
\label{sec:c02-sources}
% ============================================================================

\begin{tableaucours}{L{3.0cm}L{3.4cm}L{2.6cm}L{3.0cm}}{%
    \tableauHeader{Source} &
    \tableauHeader{Motivation} &
    \tableauHeader{Capacités} &
    \tableauHeader{Exemple}
}
Groupe criminel & Argent~: rançon, revente de données & Élevées,
organisées & Opération GAROUA \\
Employé ou complice interne & Argent, vengeance & Accès légitime &
Fraude aux fournisseurs \\
Militant (hacktiviste) & Cause politique & Variables & Défiguration de
site \\
État ou groupe soutenu & Espionnage, sabotage & Très élevées & Stuxnet
(Chapitre~\ref{chap:attaques}) \\
Opportuniste & Curiosité, défi & Faibles, outils tout faits & Balayage
automatique d'Internet \\
Erreur et accident & Aucune & Sans objet & Suppression involontaire \\
\end{tableaucours}
\vspace{-0.8em}
\begin{center}{\small\itshape Tableau~2.1. Les sources de menace}\end{center}

Une même organisation n'affronte pas toutes les sources~: une entreprise
de logistique intéresse surtout les criminels (rançon, fraude) et les
complices internes. C'est la première question d'une analyse de risques
(Chapitre~\ref{chap:risques}).

% ============================================================================
\section{Les Logiciels Malveillants}
\label{sec:c02-malware}
% ============================================================================

\begin{boxdef}
\termecle{Logiciel malveillant} (\textit{malware})~: programme conçu pour
exécuter, à l'insu de l'utilisateur, des actions nuisibles~: voler,
modifier ou détruire des données, prendre le contrôle d'une machine,
la rendre indisponible.
\end{boxdef}

\begin{tableaucours}{L{3.0cm}L{5.4cm}L{3.8cm}}{%
    \tableauHeader{Famille} &
    \tableauHeader{Caractéristique} &
    \tableauHeader{Mesure principale}
}
Virus & S'attache à un fichier et se propage quand il est ouvert &
Antivirus, prudence \\
Ver & Se propage seul par le réseau, en exploitant une faille & Correctifs,
segmentation \\
Cheval de Troie & Se présente comme un logiciel utile & Sources
logicielles maîtrisées \\
Rançongiciel & Chiffre les données et exige une rançon~; souvent précédé
d'un vol (double extorsion) & Sauvegardes hors ligne, détection \\
Logiciel espion & Collecte des informations (frappes, écran) &
Durcissement, détection \\
Porte dérobée, RAT & Donne un accès distant à l'attaquant & Filtrage
sortant, détection \\
Botnet & Réseau de machines infectées pilotées à distance & Détection
des communications de commande \\
\end{tableaucours}
\vspace{-0.8em}
\begin{center}{\small\itshape Tableau~2.2. Les familles de logiciels malveillants}\end{center}

\begin{boxalert}
\textbf{Violation typique.} Croire que l'antivirus suffit.

\smallskip
\textbf{Règle~:} l'antivirus reconnaît surtout ce qu'il a déjà vu. À
Garoua, l'attaquant a utilisé des outils d'administration légitimes, que
l'antivirus ne bloque pas. La détection d'un comportement anormal (une
connexion à 23~h~12, une copie de 40~Go) est indispensable
(Chapitre~\ref{chap:operations}).
\end{boxalert}

% ============================================================================
\section{L'Ingénierie Sociale}
\label{sec:c02-social}
% ============================================================================

\begin{boxdef}
\termecle{Ingénierie sociale}~: manipulation d'une personne pour lui faire
accomplir une action ou divulguer une information, en exploitant la
confiance, l'urgence, l'autorité ou la peur. L'\termecle{hameçonnage}
(\textit{phishing}) en est la forme la plus répandue.
\end{boxdef}

\begin{tableaucours}{L{3.2cm}L{9.0cm}}{%
    \tableauHeader{Variante} &
    \tableauHeader{Description}
}
Hameçonnage de masse & Même message envoyé à des milliers de personnes \\
Hameçonnage ciblé & Message personnalisé pour une personne ou un service \\
Fraude au président & Un faux dirigeant ordonne un virement urgent \\
Hameçonnage par SMS ou appel & Faux message d'opérateur ou de banque,
fréquent avec le Mobile Money \\
Prétexte & Un faux technicien demande un accès « pour une maintenance » \\
\end{tableaucours}
\vspace{-0.8em}
\begin{center}{\small\itshape Tableau~2.3. Les variantes de l'ingénierie sociale}\end{center}

\begin{boxscenario}{Opération GAROUA. Le courriel du 19 février}
\textbf{Constatation} (message retrouvé dans la boîte du technicien)~:
l'expéditeur affiche le nom du fournisseur de messagerie, mais le domaine
de l'adresse diffère d'une lettre~; le message annonce la suspension du
compte « dans 24~heures »~; le lien mène à une page qui reproduit la
page de connexion.

\medskip
\textbf{Analyse} (confiance élevée). Trois signaux classiques~: une
urgence artificielle, une adresse imitée, une page de connexion hors du
domaine attendu. Une formation régulière et un second facteur
d'authentification auraient rendu le vol du mot de passe inutile.
\end{boxscenario}

% ============================================================================
\section{Les Attaques contre les Réseaux et les Services}
\label{sec:c02-reseau}
% ============================================================================

\begin{tableaucours}{L{3.4cm}L{5.2cm}L{3.6cm}}{%
    \tableauHeader{Attaque} &
    \tableauHeader{Principe} &
    \tableauHeader{Propriété visée}
}
Écoute (\textit{sniffing}) & Capturer le trafic en clair & Confidentialité \\
Homme du milieu & S'interposer entre deux parties & Confidentialité,
intégrité \\
Usurpation (\textit{spoofing}) & Se faire passer pour une autre machine &
Authenticité \\
Déni de service distribué & Saturer un service depuis de nombreuses
machines & Disponibilité \\
Force brute, bourrage d'identifiants & Essayer des mots de passe, ou
réutiliser ceux qui ont fuité ailleurs & Confidentialité \\
Injection & Faire exécuter des commandes par une application
(Chapitre~\ref{chap:applications}) & Toutes \\
\end{tableaucours}
\vspace{-0.8em}
\begin{center}{\small\itshape Tableau~2.4. Attaques contre les réseaux et les
services}\end{center}

Le détail de ces attaques et des contre-mesures figure aux
Chapitres~\ref{chap:reseaux1} et~\ref{chap:reseaux2}.

% ============================================================================
\section{La Chaîne d'Attaque et MITRE ATT\&CK}
\label{sec:c02-chaine}
% ============================================================================

En 2011, des chercheurs de Lockheed Martin ont décrit l'intrusion comme une
\termecle{chaîne d'attaque} en sept étapes, chacune pouvant être
interrompue~\cite{hutchins2011}. La base MITRE ATT\&CK décrit plus
finement les \emph{tactiques} (le but de l'attaquant à une étape) et les
\emph{techniques} (le moyen), observées dans des attaques
réelles~\cite{strom2018mitre}.

\begin{tableaucours}{L{3.6cm}L{4.6cm}L{4.0cm}}{%
    \tableauHeader{Tactique ATT\&CK} &
    \tableauHeader{Opération GAROUA} &
    \tableauHeader{Mesure qui coupait la chaîne}
}
Accès initial & Mot de passe volé, accès distant & Double
authentification \\
Découverte & Exploration du réseau interne & Segmentation, détection \\
Élévation de privilèges & Faille du serveur de fichiers & Correctif \\
Exfiltration & Copie de 40~Go vers un stockage externe & Filtrage sortant,
alerte sur volume \\
Impact & Suppression des sauvegardes, chiffrement & Sauvegardes hors ligne \\
\end{tableaucours}
\vspace{-0.8em}
\begin{center}{\small\itshape Tableau~2.5. L'Opération GAROUA lue avec ATT\&CK}\end{center}

\begin{boxinfo}
\textbf{L'avantage du défenseur.} La chaîne est un inventaire
d'occasions~: il suffit de la couper à un endroit pour arrêter l'attaque,
et de voir un maillon pour la détecter. C'est le principe de la défense en
profondeur (Chapitre~\ref{chap:architecture}).
\end{boxinfo}

% ============================================================================
\section{Vulnérabilités : CVE et CVSS}
\label{sec:c02-cve}
% ============================================================================

\begin{boxdef}
\termecle{CVE} (\textit{Common Vulnerabilities and Exposures})~:
identifiant public d'une vulnérabilité connue, de la forme
CVE-année-numéro. \termecle{CVSS} (\textit{Common Vulnerability Scoring
System})~: méthode de notation de la gravité, de 0 à~10, fondée sur la
manière dont la vulnérabilité s'exploite et sur ses
conséquences~\cite{first2019cvss}.
\end{boxdef}

\begin{tableaucours}{L{3.6cm}L{8.6cm}}{%
    \tableauHeader{Métrique de base} &
    \tableauHeader{Question}
}
Vecteur d'attaque & Faut-il être sur le réseau, sur place, ou peut-on
attaquer par Internet~? \\
Complexité & L'attaque demande-t-elle des conditions particulières~? \\
Privilèges requis & Faut-il déjà un compte~? \\
Interaction & Faut-il qu'un utilisateur clique~? \\
Impacts C, I, D & Que perd-on en confidentialité, intégrité,
disponibilité~? \\
\end{tableaucours}
\vspace{-0.8em}
\begin{center}{\small\itshape Tableau~2.6. Les métriques de base du CVSS}\end{center}

\begin{boxalert}
\textbf{Violation typique.} Classer les correctifs à installer par le seul
score CVSS.

\smallskip
\textbf{Règle~:} le score mesure la gravité technique, pas le risque pour
l'organisation. Une vulnérabilité notée 7,5 sur un serveur exposé à
Internet et qui porte la facturation passe avant une vulnérabilité notée
9,0 sur un poste isolé. On croise le score avec l'exposition et la valeur
de l'actif (Chapitre~\ref{chap:systemes}).
\end{boxalert}

\begin{boxmetier}{Analyste SOC}
\textbf{Penser en chaîne.} L'analyste associe chaque alerte à une
tactique d'ATT\&CK~: une connexion distante inhabituelle relève de
l'accès initial, un nouveau compte administrateur de l'élévation de
privilèges. Deux alertes de tactiques successives sur la même machine
valent bien plus que deux alertes isolées.
\end{boxmetier}

\begin{boxmetier}{Administrateur systèmes et réseaux}
\textbf{Suivre les CVE.} L'administrateur s'abonne aux avis de sécurité de
ses éditeurs, tient l'inventaire de ses versions, et sait dire en une heure
si une nouvelle CVE critique le concerne.
\end{boxmetier}

\separateur

\begin{boxresume}
\textbf{Ce qu'il faut retenir de ce chapitre~:}
\begin{itemize}
    \item \textbf{Les sources de menace} diffèrent par leur motivation et
          leurs capacités~; on se protège d'abord contre celles qui visent
          son activité.
    \item \textbf{Les logiciels malveillants} forment des familles~;
          l'antivirus ne suffit pas contre les outils légitimes détournés.
    \item \textbf{L'ingénierie sociale} exploite la confiance et l'urgence~;
          un second facteur neutralise le vol de mot de passe.
    \item \textbf{Une attaque est une chaîne}~; ATT\&CK en décrit les
          tactiques et les techniques~; chaque maillon est une occasion.
    \item \textbf{CVE} identifie une vulnérabilité, \textbf{CVSS} en note
          la gravité technique, qui n'est pas le risque.
\end{itemize}
\end{boxresume}

\begin{boxinfo}
\textbf{Points de contrôle de l'élève.} Je sais~: classer une source de
menace~; distinguer ver, virus, cheval de Troie et rançongiciel~; repérer
trois signaux d'hameçonnage~; décrire une attaque avec les tactiques
d'ATT\&CK~; lire les métriques d'un score CVSS.
\end{boxinfo}

\bigskip

\begin{boxexo}[title={\ding{45}\quad Exercice~2.1. Quelle famille~? \niveaubase}]
Classez~: (a) un programme qui chiffre les fichiers et affiche une
demande de paiement~; (b) un programme qui se propage seul d'un serveur à
l'autre~; (c) un faux utilitaire de compression qui ouvre un accès
distant~; (d) dix mille téléphones qui envoient ensemble des requêtes.
\end{boxexo}

\begin{boxexo}[title={\ding{45}\quad Exercice~2.2. Les signaux \niveaubase}]
Listez cinq signaux qui doivent faire douter d'un courriel, et dites
lesquels étaient présents dans le courriel du 19~février.
\end{boxexo}

\begin{boxexo}[title={\ding{45}\quad Exercice~2.3. Les sources \niveaubase}]
Pour la BANQUE SANAGA, classez les sources de menace du Tableau~2.1 de la
plus à la moins préoccupante, et justifiez les deux premières.
\end{boxexo}

\begin{boxexo}[title={\ding{45}\quad Exercice~2.4. La chaîne \niveaumoyen}]
Une employée de la BANQUE SANAGA reçoit un SMS qui l'invite à
« confirmer » son compte Mobile Money~; elle saisit son code~; son compte
est vidé. Décrivez l'attaque en tactiques ATT\&CK et proposez une mesure
par tactique.
\end{boxexo}

\begin{boxexo}[title={\ding{45}\quad Exercice~2.5. CVSS \niveaumoyen}]
Une vulnérabilité exploitable par Internet, sans compte ni action de
l'utilisateur, donne le contrôle complet du serveur. Indiquez la valeur
attendue de chaque métrique de base et le niveau de gravité.
\end{boxexo}

\begin{boxexo}[title={\ding{45}\quad Exercice~2.6. Classer les correctifs \niveaumoyen}]
Trois CVE~: 9,8 sur un serveur de test isolé~; 7,5 sur le portail de
suivi exposé à Internet~; 8,1 sur les postes de la comptabilité. Dans
quel ordre les corriger~? Justifiez.
\end{boxexo}

\begin{boxexo}[title={\ding{45}\quad Exercice~2.7. Double extorsion \niveaumoyen}]
Pourquoi les sauvegardes ne suffisent-elles plus contre un rançongiciel
qui vole les données avant de les chiffrer~? Quelles mesures ajouter~?
\end{boxexo}

\begin{boxexo}[title={\ding{45}\quad Exercice~2.8. Le prétexte \niveaumoyen}]
Un homme se présente à l'agence de Bafoussam comme technicien de SAHEL
INFOGÉRANCE et demande l'accès à la salle serveur. Rédigez la procédure
que doit suivre l'agent d'accueil.
\end{boxexo}

\begin{boxexo}[title={\ding{45}\quad Exercice~2.9. Détecter plutôt que bloquer \niveauexpert}]
Pour chaque tactique du Tableau~2.5, proposez une trace que LiZbiZ aurait
pu surveiller, et la règle d'alerte correspondante.
\end{boxexo}

\begin{boxexo}[title={\ding{45}\quad Exercice~2.10. Profil de menace \niveauexpert}]
Rédigez le profil de menace d'une agence publique de paiement comme
l'ANPP~: sources, motivations, techniques probables, actifs visés, et les
trois mesures prioritaires.
\end{boxexo}

\bigskip

\begin{boxinfo}
\textbf{Pour aller plus loin.}
\begin{itemize}
    \item E.~Hutchins, M.~Cloppert et R.~Amin sur la chaîne
          d'intrusion~\cite{hutchins2011}.
    \item La base MITRE ATT\&CK~\cite{strom2018mitre}.
    \item La spécification du CVSS~v3.1~\cite{first2019cvss}.
\end{itemize}
\end{boxinfo}

\begin{boxinfo}
\textbf{Transition.} Plusieurs maillons de la chaîne se protègent par la
cryptographie~: le secret des données volées, l'intégrité des sauvegardes,
l'authenticité des connexions. Le Chapitre~\ref{chap:crypto-sym} en pose
les bases.
\end{boxinfo}
